\section{Modelos de atuadores e sensores}
\label{sec_modelos_atuadores_sensores}

Esta seção descreve as hipóteses adotadas para os modelos de atuação e de sensoriamento empregados na missão de \gls{rvdb}. 
Além disso, é mostrado como as grandezas de controle e as grandezas medidas são implementadas nas equações e nas leis de controle. 
Adotou-se a hipótese de modelos ideais, isto é, sem efeitos dissipativos como atrito, variações térmicas, elasticidade estrutural, atrasos, quantização ou ruído de medição.
Adicionalmente, assume-se comando direto de aceleração ou torque nas etapas de curta distância e final, enquanto a etapa de longa distância utiliza incrementos impulsivos de velocidade na transferência orbital.

Ainda que o hardware não seja modelado, a descrição a seguir faz referência a exemplos de atuadores e sensores empregados em missões de \gls{rvdb}, como propulsores de manobra e controle de atitude, rodas de reação, sensores inerciais, sensores estelares, câmeras e sensores ativos de distância.
Tais exemplos são incluídos apenas para contextualização física do significado das entradas e saídas consideradas no modelo ideal.

\subsection{Hipóteses gerais de idealização}
\label{subsec_hipoteses_idealizacao}

As seguintes hipóteses são adotadas durante a simulação:
\begin{itemize}
\item atuação ideal, sem dinâmicas internas de atuadores, sem atrasos e sem histerese.
\item sensoriamento ideal, com acesso direto às variáveis de estado do modelo dinâmico.
\item ausência de atrito, folgas, elasticidade estrutural e efeitos termoelásticos.
\item ausência de ruído e viés em medições de posição, velocidade, atitude e velocidade angular.
\item ausência de quantização e de limitação de resolução em sensores e atuadores.
\end{itemize}

Essas hipóteses delimitam o escopo do estudo ao desempenho lógico das leis de controle e às transições entre modelos dinâmicos, sem caracterização de desempenho de hardware ou de robustez a incertezas e perturbações não modeladas.
Em um sistema real, as variáveis utilizadas pelo controle seriam estimadas por um sistema de navegação combinando medições de múltiplos sensores e modelos dinâmicos, com filtros e tratamento de ruído e viés, conformidade \cite{wertz2011sme_smad}, o que não é considerado neste trabalho.

\subsection{Modelos de atuação}
\label{subsec_modelos_atuacao}

A seguir descreve-se como a atuação é representada em cada etapa da missão de \gls{rvdb} e como os comandos entram nas equações e nas leis de controle.

\subsubsection{Atuação translacional}
\label{subsubsec_atuacao_translacional}

Durante a aproximação de curta distância e na aproximação final, a atuação translacional é modelada como comando direto de aceleração no referencial \gls{lvlh}.
Assim, a saída das leis de controle PD translacionais é interpretada como vetor de aceleração de controle aplicado diretamente como entrada nas equações de dinâmica relativa em \gls{lvlh} (formulação \gls{hcw}).
Não há conversão de comando para o referencial inercial, nem modelagem de propulsores, empuxo, consumo de propelente ou alocação de atuadores.

Como contextualização, em uma espaçonave real essa aceleração seria tipicamente realizada por um sistema de propulsão dedicado a translação, por exemplo, propulsores químicos de baixa magnitude, propulsores elétricos em regimes adequados ou arranjos de micropropulsão, conforme requisitos de precisão, segurança e contaminação \cite{wertz2011sme_smad}.
Apesar da atuação ser ideal, impõem-se limites de projeto por meio de saturação da aceleração de comando em cada eixo (radial, along-track e cross-track), representando restrições físicas.
% Adicionalmente, a velocidade relativa admitida é limitada externamente antes de alimentar a lei de controle PD translacional, com limiar mais restritivo na região de proximidade.

\subsubsection{Atuação translacional na aproximação de longa distância}
\label{subsubsec_atuacao_translacional_longa}

Na aproximação de longa distância, a atuação é representada por incrementos impulsivos de velocidade ($\Delta v$) aplicados ao perseguidor, conforme a manobra de transferência orbital do tipo Hohmann.
Fora da aplicação do $\Delta v$, o movimento é obtido por propagação orbital do problema de dois corpos, sem acelerações de controle contínuas.

\subsubsection{Atuação de atitude}
\label{subsubsec_atuacao_atitude}

A atuação de atitude é modelada como comando direto de torque aplicado à dinâmica rotacional da espaçonave.
O torque de controle é gerado por uma lei de controle PD baseada em quaternion.
Embora utilizadas em missões reais de \gls{rvdb} \cite{wertz2011sme_smad}, não há modelagem de rodas de reação, propulsores de atitude ou magnetotorques.

Como exemplo físico, a realização desse torque em um sistema real pode ser feita por rodas de reação ou por propulsores de controle de atitude, e a escolha depende de requisitos de apontamento, saturação de momento, disponibilidade de propelente e níveis de distúrbio. No presente trabalho, tais aspectos são abstraídos pelo comando direto de torque.

Durante o modo não operacional, o torque de controle alimenta diretamente a dinâmica rotacional de corpo rígido.
Durante o modo operacional, o torque aplicado à dinâmica rotacional é modificado pela compensação do torque de reação induzido pela movimentação do \gls{mr}.

\subsubsection{Atuação do controle das juntas do MR}
\label{subsubsec_atuacao_mr}

No modo operacional do \gls{mr}, as juntas são atuadas por torque articular resultante de uma lei de controle PD em espaço articular.
A dinâmica acoplada fornece as acelerações articulares; por integração, obtêm-se as velocidades e posições articulares utilizadas na realimentação e na cinemática.
Assume-se ausência de atrito nas juntas, ausência de saturações de corrente/momento no atuador e ausência de elasticidade ou conformidade estrutural.

Em um manipulador real, os torques seriam produzidos por atuadores eletromecânicos \cite{ieee_canadarm_technical} associados a redutores e sensores de posição e velocidade, e poderiam existir limitações de corrente elétrica, saturação térmica e flexibilidade estrutural, que não são representadas aqui.

\subsection{Modelos de sensores e variáveis disponíveis ao controle}
\label{subsec_modelos_sensoriamento_variaveis}

Para realizar medições de distância entre as espaçonaves, utilizou-se sensor para esta tarefa e serão detalhadas a seguir.

\subsubsection{Estados translacionais relativos}
\label{subsubsec_sensoriamento_translacional}

As leis de controle translacionais utilizam como entradas a posição e a velocidade relativas do perseguidor em \gls{lvlh} centrado no alvo, tratadas como variáveis disponíveis de forma ideal.
Essas grandezas são obtidas do próprio modelo de dinâmica relativa em \gls{lvlh}, sem ruído ou necessidade de filtragem.
A referência translacional é especificada diretamente no referencial \gls{lvlh}, a partir da definição de uma condição geométrica desejada para viabilizar a inserção do efetuador do \gls{mr} na região de acoplamento do alvo.

Como contextualização, em missões reais a estimativa de posição e velocidade relativas pode ser obtida por sensores ópticos e ativos, como câmeras para navegação relativa e sensores de distância do tipo laser \cite{wertz2011sme_smad} rangefinder ou LiDAR, além de radar em certos perfis, combinados por filtros de estimação.
Esses elementos não são modelados, pois o estudo assume acesso direto ao estado relativo.

\subsubsection{Atitude e velocidade angular}
\label{subsubsec_sensoriamento_atitude}

A atitude do perseguidor é representada por quaternion e evolui segundo a cinemática de quaternion, integrada a partir da velocidade angular.
O quaternion é normalizado a cada passo para garantir unidade e continuidade de representação.
A lei de controle PD de atitude utiliza como entradas o quaternion atual do perseguidor e a velocidade angular do corpo, tratados como variáveis ideais fornecidas pelo modelo dinâmico correspondente (dinâmica de corpo rígido no modo não acoplado e dinâmica acoplada no modo operacional do \gls{mr}).
A atitude desejada é definida como coincidência com o \gls{lvlh} centrado no alvo.

Como referência física, em um sistema real a atitude e a velocidade angular são tipicamente obtidas por combinação de sensores inerciais (giroscópios) e sensores absolutos, como sensores estelares \cite{wertz2011sme_smad}, com fusão de dados para reduzir ruído e viés. No presente trabalho, essas grandezas são assumidas diretamente disponíveis.

\subsubsection{Estados do MR e grandezas cinemáticas}
\label{subsubsec_sensoriamento_mr_cinematica}

No modo operacional do \gls{mr}, as posições e velocidades articulares são consideradas disponíveis de forma ideal.
A cinemática direta fornece a pose do efetuador no referencial 0, e os Jacobianos linear e angular são calculados para suportar a cinemática inversa e a análise de mapeamento entre velocidades cartesianas e articulares.
A referência geométrica do efetuador é obtida a partir da posição relativa do perseguidor em \gls{lvlh}.

Como contextualização, em manipuladores reais as variáveis articulares são medidas por sensores de posição e velocidade, como encoders, e o estado do efetuador pode ser complementado por sensoriamento externo (câmeras ou sensores de distância) conforme a estratégia de acoplamento \cite{fehse2003,wertz2011sme_smad}. Esses detalhes são abstraídos pela hipótese de sensoriamento ideal.

\subsection{Representação de limites operacionais na lógica de controle}
\label{subsec_limites_operacionais_logica}

Embora não haja modelagem de hardware, limites operacionais são representados por mecanismos lógicos e saturações nas leis de controle:

\begin{itemize}
\item saturação da aceleração de comando nos canais translacionais;
\item limitação externa da velocidade relativa antes de alimentar a lei de controle PD translacional, com comutação do limite em função da distância e tempo mínimo de permanência na faixa;
\item condicionamento temporal (tempo mínimo de permanência) para comutação de modo dinâmico, reduzindo alternâncias espúrias.
\end{itemize}

\subsection{Saturações e limitadores operacionais}
\label{subsec_saturacoes_limitadores_operacionais}

As saturações e os limitadores operacionais foram introduzidos para restringir o sistema a um regime de operação compatível com os objetivos da missão e, ao mesmo tempo, reduzir a probabilidade de excitação de dinâmicas indesejadas durante eventos simultâneos, como a atuação das leis de controle, a comutação entre modelos dinâmicos e a ativação do subsistema do \gls{mr}.
Como o estudo adota modelos ideais de atuação e sensoriamento, tais mecanismos cumprem o papel de impor limites físicos e operacionais de forma explícita, evitando comandos excessivos e transientes agressivos nas variáveis de estado.

\subsubsection{Limitação da velocidade relativa na aproximação em LVLH}
\label{subsubsec_limitacao_veloc_relativa_lvlh}

Na etapa de aproximação em \gls{lvlh}, a velocidade relativa em cada eixo é limitada antes de alimentar a lei de controle PD translacional.
O limitador utiliza dois regimes em função da distância relativa: para distâncias acima de $20$ m, adota-se o limite de $\pm 0{,}30$ m/s; para distâncias iguais ou inferiores a 20 m, adota-se o limite de $\pm 0{,}03$ m/s \cite{vedant2025longrange}. 
Essa escolha visa preservar a estabilidade numérica e operacional quando múltiplos subsistemas estão ativos e quando a comutação de dinâmica se torna iminente.

A troca entre os regimes de limite de velocidade não ocorre por uma comparação instantânea de distância, mas por uma lógica de permanência temporal.
O sistema só comuta para o regime de proximidade após permanecer dentro do limiar de $1$ m por $15$ s.
Esse condicionamento reduz comutações indevidas por conta de oscilações em torno do limiar e melhora o comportamento do sistema no entorno da região de acoplamento.

\subsubsection{Saturação da aceleração de comando na lei de controle PD translacional}
\label{subsubsec_saturacao_acel_pd_translacional}

Além do limite de velocidade relativa, a lei de controle PD translacional aplica saturação na aceleração de comando em cada eixo.
Após o cálculo do comando a partir dos erros de posição e velocidade, o valor é limitado em módulo por uma aceleração máxima definida por parâmetro.
Essa saturação impede que erros transitórios produzam acelerações não realistas e reduz a possibilidade de instabilidade numérica na integração das equações \gls{hcw} quando o sistema está submetido a mudanças rápidas de regime, como a transição para a dinâmica acoplada e a ativação do subsistema do \gls{mr}.
Na implementação, a limitação de velocidade ocorre antes do cálculo do comando, e a saturação em aceleração ocorre após o cálculo do comando.

\subsubsection{Saturações no modo acoplado com o \gls{mr}}
\label{subsubsec_saturacoes_modo_acoplado_mr}

No modo acoplado, são definidos limites adicionais associados às variáveis do subsistema dinâmico da espaçonave e do \gls{mr}.
Em particular, são estabelecidos limites de aceleração, velocidade e posição para as juntas, bem como limites de aceleração e velocidade para a parte da espaçonave no modelo acoplado.
Esses limites atuam como proteções operacionais no momento em que o sistema passa a considerar o acoplamento dinâmico e em que a cinemática inversa, a cinemática direta e os Jacobianos passam a operar em conjunto com as leis de controle.
Adicionalmente, o torque articular comandado pela lei de controle PD do \gls{mr} é limitado por um valor máximo por junta, representando restrição operacional de atuação.

\subsubsection{Motivação física para os limitadores}
\label{subsubsec_motivacao_fisica_limitadores}

A motivação principal para os limitadores é física: garantir que a dinâmica relativa e a dinâmica acoplada permaneçam em um regime coerente com a aproximação, minimizando transientes que poderiam comprometer a estabilidade do sistema quando várias ações ocorrem simultaneamente.
Na aproximação no referencial \gls{lvlh}, a limitação de velocidade ajuda a manter a estabilidade do sistema, favorecendo uma entrada suave na janela geométrica de comutação.
A saturação de aceleração restringe a agressividade da lei de controle PD translacional diante de erros momentâneos, evitando picos de comando.
No modo operacional, as saturações associadas ao \gls{mr} e à espaçonave reduzem a chance de respostas não físicas em presença de acoplamentos e contribuem para a estabilidade numérica e para a coerência do regime de operação durante a aproximação final.